\chapter{Các mục chờ bổ sung và khuyến nghị}
\label{app:cho-bo-sung}

Phụ lục này liệt kê những phần chưa có bằng chứng đầy đủ, để người đọc biết chính xác phạm vi
của báo cáo và quyết định phần nào cần bổ sung. Mỗi mục ghi: câu hỏi, trạng thái, lý do còn
thiếu, công sức ước tính, tác động lên kết luận, và khuyến nghị. Không có mục nào được coi là
``đã trả lời'' nếu chỉ có cấu hình mà chưa có log.

\section{Bảng trạng thái}
\begin{small}
\begin{longtable}{@{}p{.3\textwidth} p{.16\textwidth} p{.1\textwidth} p{.12\textwidth} p{.2\textwidth}@{}}
\toprule
\textbf{Hạng mục} & \textbf{Trạng thái} & \textbf{Công sức} & \textbf{Tác động} & \textbf{Khuyến nghị} \\
\midrule
Giải chính chế độ A (cùng thời gian) & Chưa có & 4--8 giờ & Cao & Nên bổ sung \\
Lặp nhiều seed & Chưa có & 4--8 giờ & Cao & Nên bổ sung \\
Tập khai cuộc giữ riêng & Chưa có & 1--2 giờ & Cao & Nên bổ sung \\
A6 tăng cường đối xứng & Chưa có checkpoint & 1--2 giờ & Trung bình & Tuỳ chọn \\
A7b đối đầu trọng số & Chưa chạy & 30 phút & Trung bình & Nên bổ sung \\
A8b root tactics & Chưa chạy & 1--2 giờ & Trung bình & Nên bổ sung \\
A5b khởi tạo RL & Chưa có cặp đủ & 8--12 giờ & Trung bình & Cân nhắc \\
Top-1/top-3 trên tập giữ riêng & Đã có & — & — & Đã xong \\
Calibration value theo pha & Đã có & — & — & Đã xong \\
Phân tích ván sâu & 20 ván & 2--3 giờ & Trung bình & Có thể thêm \\
PUCT Python vs Cython & Số cũ & 30 phút & Thấp & Tuỳ chọn \\
Gom suy luận GPU & Chưa đo & 1--2 giờ & Thấp & Tuỳ chọn \\
Đối chiếu engine ngoài & Chưa có & 2--4 giờ & Cao & Khó, cân nhắc \\
Biến thể luật/bàn & Chưa có & 4--8 giờ & Cao & Ngoài phạm vi \\
Khảo sát người dùng UI & Chưa có & 2--3 giờ & Thấp & Tuỳ chọn \\
\bottomrule
\end{longtable}
\end{small}

\section{Giải chính chế độ A}
\paragraph{Câu hỏi.} Với cùng 1\,000 ms mỗi nước, thứ hạng tám tác tử có khác chế độ B không?
\paragraph{Trạng thái.} Chưa có. Báo cáo chỉ có RR8 chế độ B (ngân sách cố định theo loại thuật
toán) và các phép đo A3/A4/A5 ở ngân sách cố định.
\paragraph{Lý do.} Chạy chế độ A cho tám tác tử tốn khoảng 4--8 giờ vì mỗi nước chờ đủ 1 giây,
kể cả khi tác tử đã tìm xong. Với 2\,800 ván và khoảng 30 nước/ván, tổng thời gian chờ rất lớn.
\paragraph{Tác động.} Cao. Đây là phần còn thiếu quan trọng nhất để trả lời RQ1 đầy đủ. Chế độ
A có thể đảo thứ hạng giữa Alpha-Beta và PUCT/MCTS vì Alpha-Beta dùng thời gian hiệu quả hơn
trên mỗi nút, còn PUCT có thể dùng nhiều mô phỏng hơn.
\paragraph{Khuyến nghị.} Nên bổ sung, ít nhất cho các cặp L3--L4, L3--L7, L3--L8, L7--L8, thay vì
toàn bộ 28 cặp. Chạy 20 khai cuộc × 2 màu × 5 cặp = 200 ván, khoảng 1--2 giờ.

\section{Lặp nhiều seed}
\paragraph{Câu hỏi.} Kết quả có ổn định qua các seed gốc khác nhau không?
\paragraph{Trạng thái.} Chưa có. Mọi run dùng seed 2026 (hoặc 42 cho A5).
\paragraph{Lý do.} Mỗi lần chạy lại tốn thời gian; L8 đặc biệt tốn vì tự đấu.
\paragraph{Tác động.} Cao. Một chênh lệch 5 điểm phần trăm có thể do seed, không do phương pháp.
Hiện báo cáo chỉ có một quan sát cho mỗi cấu hình, nên không thể tách nhiễu seed.
\paragraph{Khuyến nghị.} Nên bổ sung cho các cặp chính, không cần cho cả tám tác tử. Chạy lại
RR8 với 3 seed (2026, 2027, 2028) trên 20 khai cuộc, hoặc chạy L8 với 3 seed nếu có thời gian.

\section{Tập khai cuộc giữ riêng}
\paragraph{Câu hỏi.} Kết quả có giữ nguyên trên các khai cuộc chưa dùng trong phát triển/gating?
\paragraph{Trạng thái.} Chưa có. Tập 100 khai cuộc được dùng cho cả phát triển, A3/A4, A5, RR8
và gating L8.
\paragraph{Lý do.} Cần sinh thêm khai cuộc và chạy lại các cặp chính.
\paragraph{Tác động.} Cao. Đây là cách chuẩn để tránh thiên lệch chọn mẫu. Đặc biệt quan trọng
cho L8, vì checkpoint được chọn bằng gating trên cùng tập khai cuộc dùng để đánh giá cuối.
\paragraph{Khuyến nghị.} Nên bổ sung. Sinh 50 khai cuộc mới với seed khác, chạy L7--L8, L8--L3,
L3--L7, mỗi cặp 100 ván. Khoảng 1--2 giờ.

\section{A6: tăng cường đối xứng}
\paragraph{Câu hỏi.} Huấn luyện có tăng cường D4 có cải thiện sức chơi so với không tăng cường?
\paragraph{Trạng thái.} Chưa có cặp checkpoint. Có cấu hình \texttt{experiments/a6\_augment.yaml}
nhưng chưa chạy hai lần.
\paragraph{Lý do.} Cần huấn luyện hai mô hình cùng dữ liệu, cùng epoch, khác cờ augment. Mỗi
lần khoảng 20--40 phút trên GPU.
\paragraph{Tác động.} Trung bình. Tăng cường D4 thường giúp ích nhưng mức độ phụ thuộc dữ liệu;
không phải kết luận trung tâm của báo cáo.
\paragraph{Khuyến nghị.} Tuỳ chọn. Nếu làm, dùng 15 epoch thay vì 30 để tiết kiệm, và đấu trên
20 khai cuộc.

\section{A7b: đối đầu hai bảng trọng số}
\paragraph{Câu hỏi.} Bảng trọng số mặc định có chơi mạnh hơn bảng bỏ phòng thủ không?
\paragraph{Trạng thái.} Chưa chạy. Có cấu hình \texttt{experiments/a7b\_head\_to\_head.yaml}.
A7 cũ chỉ so giá trị gốc trên 10 thế, cho thấy 9/10 thế khác giá trị.
\paragraph{Lý do.} Cần chạy 100 ván; rất nhanh nhưng chưa được ưu tiên.
\paragraph{Tác động.} Trung bình. Trả lời trực tiếp câu hỏi ``thành phần phòng thủ có cần không''
ở mức sức chơi, không chỉ giá trị gốc.
\paragraph{Khuyến nghị.} Nên bổ sung; chỉ mất khoảng 30 phút.

\section{A8b: chiến thuật tại gốc}
\paragraph{Câu hỏi.} Bật \texttt{root\_tactics} cho MCTS có cải thiện điểm trước L2/L3 không?
\paragraph{Trạng thái.} Chưa chạy. Có cấu hình \texttt{experiments/a8b\_root\_tactics.yaml}.
\paragraph{Lý do.} MCTS rollout nhẹ tốn khoảng 0,6 giây/nước; 600 ván mất 1--2 giờ.
\paragraph{Tác động.} Trung bình. Kết quả A3/A4 cho thấy MCTS bỏ chặn rất nhiều, nên giả thuyết
là bật chiến thuật sẽ cải thiện mạnh. Nếu đúng, đây là bằng chứng ``sửa lỗi bằng luật'' rẻ hơn
tăng mô phỏng.
\paragraph{Khuyến nghị.} Nên bổ sung nếu còn thời gian; đây là thí nghiệm có tỷ lệ tác động/công
sức tốt.

\section{A5b: khởi tạo RL}
\paragraph{Câu hỏi.} RL khởi tạo từ SL có tốt hơn RL khởi tạo ngẫu nhiên không?
\paragraph{Trạng thái.} Chưa có cặp đủ số ván. Có cấu hình \texttt{experiments/a5b\_rl\_init.yaml}
và các run smoke, nhưng chưa chạy 50k episode cho cả hai nhánh.
\paragraph{Lý do.} Mỗi nhánh 50k episode tốn nhiều giờ; chạy hai nhánh tốn gấp đôi.
\paragraph{Tác động.} Trung bình. Trả lời câu hỏi về vai trò của tri thức chuyên gia trong RL.
\paragraph{Khuyến nghị.} Cân nhắc; chỉ nên làm nếu có GPU mạnh và thời gian.

\section{Những phần đã hoàn thành trong bản này}
Để tránh hiểu nhầm rằng mọi thứ đều thiếu, bảng dưới liệt kê các phần đã có bằng chứng mới:
\begin{itemize}
  \item RR8: 2\,800 ván, tám tác tử, 50 khai cuộc đổi màu.
  \item E3: 800 ván, tám cấp tự đấu, 50 khai cuộc đổi màu.
  \item Top-1/top-3 và calibration value trên 29\,161 mẫu test.
  \item Phân tích chiến thuật trên hàng nghìn ván, đếm bỏ thắng/bỏ chặn.
  \item Mười ván phân tích sâu với bàn TikZ và cơ chế kết thúc.
  \item A1/A2/A7 benchmark và A3/A4/A5 (dữ liệu lịch sử, có đối chiếu).
\end{itemize}

\section{Khuyến nghị ưu tiên}
Nếu chỉ có 4--6 giờ để bổ sung, thứ tự đề xuất là:
\begin{enumerate}
  \item Tập khai cuộc giữ riêng (1--2 giờ) — cải thiện giá trị khoa học lớn nhất trên mỗi giờ.
  \item A7b và A8b (1,5--2,5 giờ) — trả lời hai câu hỏi ablation còn trống.
  \item Chế độ A cho 5 cặp chính (1--2 giờ) — trả lời RQ1 đầy đủ hơn.
  \item Lặp 3 seed cho RR8 trên 20 khai cuộc (2--4 giờ) — đo độ ổn định.
\end{enumerate}
Các mục còn lại có thể để hướng phát triển. Việc chọn hay bỏ do người đọc quyết định; phụ lục
này chỉ nêu điều kiện và hệ quả.